% الجزء: sets-functions-relations
% الفصل: relations
% القسم: equivalence-relations

\documentclass[../../../include/open-logic-section]{subfiles}

\begin{document}

\olfileid[ar]{sfr}{rel}{eqv}

\olsection{علاقات التكافؤ}

تجتمع في علاقة الهوية على المجموعة ثلاث خصائص: الانعكاسية والتناظر
والتعدية. والعلاقات~$R$ الجامعة لهذه الثلاث كثيرة الورود.

\begin{defn}[علاقة التكافؤ]
\emph{علاقة التكافؤ} علاقة $R \subseteq A^2$ تجمع الانعكاسية والتناظر
والتعدية. ويقال إن !!^{element}s $x$ و$y$ من~$A$
\emph{متكافئان وفق $R$} متى كان~$Rxy$.
\end{defn}

ومن علاقة التكافؤ تؤخذ \emph{فئات التكافؤ}. وذلك أن العلاقة تقسم
المجال كتلًا: تتعلق الأشياء في كل كتلة بعضها ببعض، ولا يتعلق شيء
من كتلة بشيء من أخرى. وقد ينفع أن يكون الكلام على هذه الكتل
\emph{مباشرة}؛ فلهذا نعرّفها كما يأتي.

\begin{defn}\ollabel{def:equivalenceclass}
لتكن $R \subseteq A^2$ علاقة تكافؤ. إذا أخذ $x \in A$ كانت
\emph{فئة تكافؤ} $x$ في~$A$ هي المجموعة $\equivrep{x}{R}
= \Setabs{y \in A}{Rxy}$. وأما \emph{مجموعة قسمة} $A$ بالنسبة
إلى~$R$ فهي $\equivclass{A}{R} = \Setabs{\equivrep{x}{R}}{x \in A}$؛
أي المجموعة الجامعة لفئات التكافؤ هذه.
\end{defn}

وتشهد النتيجة الآتية لصحة هذا التعريف، إذ تبين أن فئات التكافؤ
هي حقًا الكتل التي تتجزأ إليها~$A$:

\begin{prop}
لتكن $R \subseteq A^2$ علاقة تكافؤ. حينئذ يتكافأ $Rxy$ مع
$\equivrep{x}{R} = \equivrep{y}{R}$.
\end{prop}

\begin{proof}
أما من اليسار إلى اليمين، فنفرض $Rxy$ ونأخذ $z \in
\equivrep{x}{R}$. فبالتعريف $Rxz$، وبما أن $R$ تكافؤ يلزم $Ryz$.
وتفصيل اللزوم: من $Rxy$ وتناظر~$R$ نحصل على $Ryx$، ثم من
$Rxz$ وتعدية~$R$ نحصل على~$Ryz$. إذن $z \in \equivrep{y}{R}$،
وبعموم الاختيار $\equivrep{x}{R} \subseteq \equivrep{y}{R}$.
وبمثل ذلك $\equivrep{y}{R} \subseteq \equivrep{x}{R}$؛ فالامتدادية
تقضي بأن $\equivrep{x}{R} = \equivrep{y}{R}$.

وأما من اليمين إلى اليسار، فنفرض $\equivrep{x}{R} =
\equivrep{y}{R}$. وانعكاسية $R$ تعطي $Ryy$، ومنه
$y \in \equivrep{y}{R}$. فيكون $y \in \equivrep{x}{R}$ أيضًا،
لأن $\equivrep{x}{R} = \equivrep{y}{R}$ بالفرض؛ أي $Rxy$.
\end{proof}

\begin{ex}
ومن أمثلة التكافؤ الحساب بترديد عدد. خذ $a$ و$b$ و$n \in \PosInt$.
نقول $a \equiv_n b$ إذا وفقط إذا اتفق باقي قسمة $a$ على~$n$
وباقي قسمة $b$ على~$n$. وصورته الرمزية أن $a \equiv_n b$ إذا وفقط
إذا وجد $k \in \Int$ يحقق $a - b = kn$. فـ$\equiv_n$ علاقة تكافؤ
لكل~$n$، وعدد الفئات المتمايزة $n$ بالضبط، وهي الناشئة عن~$\equiv_n$؛ أي إن
$\equivclass{\Nat}{\equiv_n}$ عدد !!{element}sها $n$.
فمنها الأعداد القابلة للقسمة على $n$ بلا باق، وهي
$\equivrep{0}{\equiv_n}$؛ ثم ما بقي من قسمته على $n$ العدد~$1$،
وهو $\equivrep{1}{\equiv_n}$؛ و\ldots؛ إلى ما بقي من قسمته على~$n$
العدد~$n-1$، وهو~$\equivrep{n-1}{\equiv_n}$.
\end{ex}

\begin{prob}
برهن على أن $\equiv_n$ علاقة تكافؤ لكل $n \in \PosInt$، وأن
عدد عناصر $\equivclass{\Nat}{\equiv_n}$ هو $n$ بالضبط.
\end{prob}

\end{document}
